\subsection{Kreuzprodukt in Koordinatendarstellung}
\label{subsec:vector_algebra_cross_product_coordinates}

Wir stellen $\vec a,\vec b\in\R^3$ mithilfe der kanonischen Einheitsvektoren
dar:
\[
  \vec a=a_1\vec e_1+a_2\vec e_2+a_3\vec e_3,
  \qquad
  \vec b=b_1\vec e_1+b_2\vec e_2+b_3\vec e_3.
\]
Mit den Kreuzprodukten der kanonischen Einheitsvektoren folgt durch
Ausmultiplizieren
\begin{ClozeArea}[34]{16}
  \[
    \begin{aligned}
      \vec a\times\vec b
      &=(a_1\vec e_1+a_2\vec e_2+a_3\vec e_3)
        \times(b_1\vec e_1+b_2\vec e_2+b_3\vec e_3)\\[1.5em]
      &=a_1b_1\overbrace{(\vec e_1\times\vec e_1)}^{=\,\vec 0}
        +a_1b_2\overbrace{(\vec e_1\times\vec e_2)}^{=\,\vec e_3}
        +a_1b_3\overbrace{(\vec e_1\times\vec e_3)}^{=\,-\vec e_2}
        \\[0.9em]
      &\quad
        +a_2b_1\overbrace{(\vec e_2\times\vec e_1)}^{=\,-\vec e_3}
        +a_2b_2\overbrace{(\vec e_2\times\vec e_2)}^{=\,\vec 0}
        +a_2b_3\overbrace{(\vec e_2\times\vec e_3)}^{=\,\vec e_1}
        \\[0.9em]
      &\quad
        +a_3b_1\overbrace{(\vec e_3\times\vec e_1)}^{=\,\vec e_2}
        +a_3b_2\overbrace{(\vec e_3\times\vec e_2)}^{=\,-\vec e_1}
        +a_3b_3\overbrace{(\vec e_3\times\vec e_3)}^{=\,\vec 0}
        \\[1.3em]
      &=(a_2b_3-a_3b_2)\vec e_1
        +(a_3b_1-a_1b_3)\vec e_2
        +(a_1b_2-a_2b_1)\vec e_3\\[1em]
      &=\begin{pmatrix}
          a_2b_3-a_3b_2\\
          a_3b_1-a_1b_3\\
          a_1b_2-a_2b_1
        \end{pmatrix}.
    \end{aligned}
  \]
\end{ClozeArea}

\begin{KeyPointBox}{Koordinatenformel}{vector_algebra_cross_product_coordinate_formula}
  Für $\vec a,\vec b\in\R^3$ gilt
  \[
    \vec a\times\vec b
    =\begin{ClozeContinuation}[22]{3}
      \begin{pmatrix}
        a_2b_3-a_3b_2\\
        a_3b_1-a_1b_3\\
        a_1b_2-a_2b_1
      \end{pmatrix}
    \end{ClozeContinuation}.
  \]
\end{KeyPointBox}

Aus der Koordinatenformel in
Merke~\ref{key:vector_algebra_cross_product_coordinate_formula} gewinnen wir
die Determinantenschreibweise schrittweise. Zunächst zerlegen wir den
Koordinatenvektor in die kanonischen Einheitsvektoren:
\begin{ClozeArea}[34]{12}
  \[
    \begin{aligned}
      \vec a\times\vec b
      &=\begin{pmatrix}
          a_2b_3-a_3b_2\\
          a_3b_1-a_1b_3\\
          a_1b_2-a_2b_1
        \end{pmatrix}
        =(a_2b_3-a_3b_2)\vec e_1
        +(a_3b_1-a_1b_3)\vec e_2
        +(a_1b_2-a_2b_1)\vec e_3
        \\[1.4em]
      &=\det\!\begin{pmatrix}
          a_2&b_2\\
          a_3&b_3
        \end{pmatrix}\vec e_1
        -\det\!\begin{pmatrix}
          a_1&b_1\\
          a_3&b_3
        \end{pmatrix}\vec e_2
        +\det\!\begin{pmatrix}
          a_1&b_1\\
          a_2&b_2
        \end{pmatrix}\vec e_3
        \\[1.4em]
      &=\det\!\begin{pmatrix}
          \MathHighlight{\vec e_1}&a_1&b_1\\
          \MathHighlight{\vec e_2}&a_2&b_2\\
          \MathHighlight{\vec e_3}&a_3&b_3
        \end{pmatrix}.
    \end{aligned}
  \]
\end{ClozeArea}
Die letzte Gleichheit liest die Entwicklung nach der gelb markierten ersten
Spalte rückwärts. Die Matrix ist gegenüber der üblichen Zeilendarstellung
transponiert, sodass die Koordinaten von $\vec a$ und $\vec b$ in den Spalten
stehen. Wegen $\det(M^{\mathsf T})=\det(M)$ ändert die Transposition den
Determinantenwert nicht. Diese Schreibweise dient nur als formale Merkhilfe;
definiert wurde das Kreuzprodukt geometrisch.

\begin{KeyPointBox}{Determinantenschreibweise}{vector_algebra_cross_product_determinant_notation}
  Das Kreuzprodukt lässt sich formal schreiben als
  \[
    \vec a\times\vec b
    =\begin{ClozeContinuation}[22]{4}
      \det\!\begin{pmatrix}
        \vec e_1&a_1&b_1 \\
        \vec e_2&a_2&b_2 \\
        \vec e_3&a_3&b_3
      \end{pmatrix}
    \end{ClozeContinuation}.
  \]
\end{KeyPointBox}

\begin{ExampleBox}{Kreuzprodukt aus Koordinaten}{vector_algebra_cross_product_coordinates_example}
  Für
  \[
    \vec a=\begin{pmatrix}2\\1\\-2\end{pmatrix},
    \qquad
    \vec b=\begin{pmatrix}-1\\2\\-1\end{pmatrix}
  \]
  ergibt sich
  \[
    \vec c=\vec a\times\vec b
    =\begin{ClozeContinuation}[22]{4}
      \begin{pmatrix}
        1\cdot(-1)-(-2)\cdot2\\
        (-2)\cdot(-1)-2\cdot(-1)\\
        2\cdot2-1\cdot(-1)
      \end{pmatrix}
      =\begin{pmatrix}3\\4\\5\end{pmatrix}
    \end{ClozeContinuation}.
  \]
  Die Orthogonalität zu beiden Ausgangsvektoren wird durch das Skalarprodukt
  bestätigt:
  \[
    \vec c\cdot\vec a=\Cloze{5cm}{6+4-10=0},
    \qquad
    \vec c\cdot\vec b=\Cloze{5cm}{-3+8-5=0}.
  \]
  Der aufgespannte Flächeninhalt beträgt
  \[
    A=|\vec a\times\vec b|
    =\Cloze{7cm}{\sqrt{3^2+4^2+5^2}=\sqrt{50}}.
  \]
\end{ExampleBox}
